\section{Conclusion}
\label{sec:conclusion}

A generative language model defines a measure on token strings; classes are reached only through
a parsing map that leaks probability, and the class distribution can only be estimated. A model
that returns a probability for every option defines the measure on $(C, 2^C)$ directly. On four
datasets that was enough to apply calibration analysis, threshold selection, escalation,
importance-sampled evaluation and perturbation-based explanations to Jev with standard tools,
and to score it with the same code as supervised classifiers; obtaining the same object from a
generative model took a scoring procedure of its own and far more computation per prediction.
What the toolbox found is mixed, and that is its value: accurate, usable probabilities on text;
a bias towards the adverse class on tables that a threshold or a second model repairs in one
case and not in the other; and guided sampling and escalation that pay off only where the
reported confidence is informative. The first thing the classical toolbox reports about a class
distribution is where it can be trusted.
